\section{Do forecasts change for the right reason?}
\label{sec:evidence}

A key forecasting skill is to appropriately incorporate new evidence as it
arrives; for example, classic work on judgmental forecasting identifies repeated
belief revision as a hallmark of strong human forecasters
\citep{tetlock2005expert}. We argue that effective updating requires a forecaster
to determine: (1) which evidence is relevant in predicting an outcome and (2) the
extent to which the new evidence should affect the forecast. We assess these
capacities in two experiments. First, we pair real-world prediction-market
questions with controlled synthetic evidence, allowing us to test whether models
update selectively and in the direction the evidence warrants. Second, we use a
controlled synthetic environment in which the same evidence has different
predictive effects across contexts, allowing us to test whether models can infer
the relationship that determines how strongly the evidence should affect their
forecast.

\subsection{Exp.~1: Can models directionally update?}
\label{sec:exp1}

\noindent\textbf{Design.}
We ask nine systems to forecast 100 unresolved binary Polymarket questions
with web search, then revise each forecast after one fictional news update with
search disabled (Fig.~\ref{fig:exp1-pipeline}). Updates support YES, support NO,
or concern the same topic without bearing on the outcome; these labels are hidden
from the model. This yields 30,094 valid updates
(App.~\ref{app:exp1}).

\noindent\textbf{Finding: models track the intended packet direction.}
Every system leaves its forecast unchanged more often for topical controls.
Conditional on moving, models move $1.9$--$7.5\times$ farther for directional
than topical updates. For revisions of at least three percentage points,
directional consistency ranges from $.887$ to $.990$, versus $.688$ for a
word-only classifier evaluated on held-out markets. Human validation and
uncertainty estimates are in App.~\ref{app:exp1-threshold-robustness}.

\begin{figure}[!b]
\centering
\begin{subfigure}[c]{0.525\textwidth}
\centering
\includegraphics[width=.92\linewidth]{exp1_direction_selective_updating.pdf}
\caption{}
\label{fig:exp1-bars}
\end{subfigure}\hspace{0.01\textwidth}
\begin{subfigure}[c]{0.455\textwidth}
\centering
\footnotesize
\setlength{\tabcolsep}{3pt}
\begin{tabular*}{\linewidth}{@{}l@{\extracolsep{\fill}}rccc@{}}
\toprule
\textbf{Model} & \textbf{Size} & \textbf{Abst.} & \textbf{Sens.}$\uparrow$ & \textbf{EHC$_{\geq3}$}$\uparrow$ \\
\midrule
\multicolumn{5}{@{}l}{\emph{Hosted API (parameters undisclosed)}}\\
\micon{claude.png}~Opus~4.8 & --- & 71.4\% & \cellcolor{heatgreen}$7.0\times$ & \cellcolor{heatgreen}.990 \\
\micon{openai.png}~GPT-5.4 & --- & 8.4\% & \cellcolor{heatgreen}$7.5\times$ & \cellcolor{heatgreen}.987 \\
\micon{deepseek.png}~V4-Pro & --- & 44.4\% & \cellcolor{heatyellow}$4.9\times$ & \cellcolor{heatyellow}.977 \\
\midrule
\multicolumn{5}{@{}l}{\emph{Open-weight, descending parameters} $\downarrow$}\\
\micon{qwen.pdf}~2.5-72B & 72B & 24.2\% & \cellcolor{heatorange}$3.2\times$ & \cellcolor{heatyellow}.963 \\
\micon{llamma.pdf}~3.3-70B & 70B & 10.7\% & \cellcolor{heatyellow}$5.7\times$ & \cellcolor{heatyellow}.965 \\
\micon{llamma.pdf}~3.1-70B & 70B & 2.5\% & \cellcolor{heatyellow}$4.5\times$ & \cellcolor{heatyellow}.961 \\
\micon{qwen.pdf}~2.5-32B & 32B & 11.1\% & \cellcolor{heatorange}$2.8\times$ & \cellcolor{heatyellow}.964 \\
\micon{qwen.pdf}~2.5-14B & 14B & 26.1\% & \cellcolor{heatred}$2.2\times$ & \cellcolor{heatorange}.943 \\
\micon{llamma.pdf}~3.1-8B & 8B & 1.1\% & \cellcolor{heatred}$1.9\times$ & \cellcolor{heatred}.887 \\
\bottomrule
\end{tabular*}
\caption{}
\label{fig:exp1-table}
\label{tab:consistency}
\end{subfigure}
\caption{\textbf{Models move much more for evidence than for topical
chatter, and usually move in the implied direction.} \textbf{(a)} Average
revision after supporting, contrary, and topical updates. \textbf{(b)} The
table summarizes how often models leave topical updates alone, how much farther
they move for directional evidence conditional on moving (Sens.), and how often
a move of at least three percentage points has the right sign (EHC$_{\geq3}$). Intervals and robustness checks are in
App.~\ref{app:exp1-threshold-robustness}.}
\label{fig:exp1-updating}
\end{figure}

\noindent\textbf{Takeaway.}
Lexical cues identify original-packet relevance ($.989$ accuracy). On 20
unvalidated context-reversal families, GPT-5.6 reverses both directions in
85\% of pairs, with 0.5-point broken-link movement; the original constrained
local deployments reach 25--35\% and 9.75--13.25 points in the matched
single-repeat comparison. A matched Qwen3 follow-up reaches 60\% paired
reversal with thinking enabled (App.~\ref{app:exp1-context-reversal}). This
exploratory extension uses different inference settings and requires independent
material validation. Experiment~2
tests contextual relationship selection under a controlled generative process.

\subsection{Exp.~2: can models infer how strongly evidence should matter?}
\label{sec:exp2}

\begin{figure}[!t]
\centering
\includegraphics[
  width=.90\linewidth,
  trim=11pt 14.5pt 3.5pt 7pt,
  clip
]{exp2_coin_city_prompt_arms.pdf}
\caption{\textbf{Coin City turns contextual reasoning into a paired
comparison.} Every arm shares the query and numerical cases. The context arm
adds a description connecting City C to one displayed media environment; the
opposite-context arm changes only that sentence. Target cases then accumulate,
allowing local evidence to confirm or correct the contextual choice.}
\label{fig:coin-city-prompt}
\end{figure}

\noindent\textbf{Design.}
Experiment~1 tests which evidence warrants an update. Forecasting also requires
inferring how strongly that evidence should matter. In the opening coin-flip
example (Fig.~\ref{fig:world-knowledge-evidence}), the same observation warrants
different predictions depending on the process believed to have generated it.

We systematize this ambiguity in \emph{Coin City}: the same campaign news moves
polls strongly in one reference city and weakly in another. Context identifies the
better analogue for target City C, while zero to four City C cases provide direct
evidence that can confirm or override it. Each city's polling response follows
\begin{equation}
  \Delta p_{ij}=g_j x_{ij}+\epsilon_{ij},
  \qquad \epsilon_{ij}\sim\mathcal{N}(0,5^2),
  \label{eq:coin-city-dgp}
\end{equation}
where $x_{ij}$ is signed campaign news and $g_j$ is the city's stable sensitivity,
which is never stated: the model must infer it from examples.

We compare target cases alone, reference cases without target context, correct
context, and opposite context, holding numerical evidence fixed across the paired
arms (Fig.~\ref{fig:coin-city-prompt}). A fifth arm replaces familiar descriptions
with episode-varying labels whose meaning must be inferred from the references.
Forecast error tests whether context helps; the forecast-implied response to news
tests which reference the model selects (App.~\ref{app:coin-city-results}).

\noindent\textbf{Finding: the forecast follows the contextual clue.}
Figure~\ref{fig:coin-city-results} shows that, when City C provides few direct
observations, context guides which reference city's response is used to forecast.
The correct context lowers forecast error, while the wrong analogue increases
error. Moreover, the effect implies the model is quantifying the importance of
the contextual clue: after controlling for the underlying response regime, the
forecast-implied slope correlates $r=.51$--$.86$ with the fitted slope of the
cue-selected reference city, while its correlation with the other reference
remains near zero. As City C observations accumulate, context matters less, as
forecasts increasingly rely on direct local evidence.

\begin{figure}[!t]
\centering
\includegraphics[width=\linewidth]{exp2_coin_city_results.pdf}
\vspace{3pt}
\begin{minipage}{\linewidth}
\centering
\footnotesize
\setlength{\tabcolsep}{4pt}
\resizebox{\linewidth}{!}{%
\begin{tabular}{l r rrrr rr}
\toprule
& & \multicolumn{4}{c}{Forecast MAE [95\% CI]} & \multicolumn{2}{c}{$\rho(\hat g,g_C)$} \\
\cmidrule(lr){3-6}\cmidrule(lr){7-8}
Model & Paired $n$ & Target only & No C context & $+$ C context &
$+$ wrong C context & No context & $+$ context \\
\midrule
\micon{claude.png}~Opus~4.8 & 250 & 2.32 [2.20,\,2.46] & 2.65 [2.41,\,2.89] & \textbf{1.82} [1.64,\,2.02] & 4.20 [3.87,\,4.54] & $-.06$ & \textbf{.70} \\
\micon{claude.png}~Opus~5 & 250 & 2.36 [2.12,\,2.59] & 2.62 [2.38,\,2.84] & \textbf{1.76} [1.57,\,1.99] & 4.25 [3.92,\,4.57] & $-.01$ & \textbf{.62} \\
\micon{openai.png}~GPT-5.6 & 250 & 3.60 [3.28,\,3.94] & 2.52 [2.33,\,2.75] & \textbf{1.60} [1.42,\,1.77] & 4.23 [3.90,\,4.55] & $-.04$ & \textbf{.68} \\
\micon{deepseek.png}~V4-Pro & 250 & \textbf{2.45} [2.27,\,2.61] & 2.84 [2.55,\,3.15] & 2.47 [2.19,\,2.77] & 4.36 [4.01,\,4.74] & $-.03$ & \textbf{.52} \\
\micon{kimi.png}~K3 & 250 & 3.35 [3.04,\,3.67] & 2.55 [2.32,\,2.78] & \textbf{1.89} [1.66,\,2.13] & 4.31 [3.98,\,4.62] & $-.02$ & \textbf{.59} \\
\micon{gemini.png}~3.6 Flash & 250 & 2.53 [2.35,\,2.73] & 2.55 [2.31,\,2.80] & \textbf{2.17} [1.89,\,2.48] & 4.29 [3.94,\,4.66] & $.08$ & \textbf{.56} \\
\bottomrule
\end{tabular}
}
\end{minipage}
\caption{\textbf{Context improves forecasts before target cases and helps assign
the response pattern across six systems.} Giving the correct analogue lowers
error before City C has a history; naming the wrong analogue raises it. As direct
City C cases accumulate, forecasts rely less on either contextual cue. Full
results by model are in App.~\ref{app:coin-city-results}.}
\label{fig:coin-city-results}
\end{figure}

\subsubsection{Infer a convention without stable semantics}
\label{sec:novel-mapping}

However, a different explanation for the relational mapping above is that the
model may already associate familiar cues, like ``national news,'' with a
stronger response and ``local news'' with a weaker one, rather than actually
inferring the mapping from the reference cases. We remove this shortcut by
replacing these cues with the arbitrary labels KIV and ZOR. Their meanings are
randomized across episodes, so neither label has any stable association with
strong or weak responsiveness. Twelve of fifteen tested systems recover the episode-local mapping before any
target cases, including all tested local checkpoints above 8B
(App.~\ref{app:coin-city-symbol-context}), although recovery does not always
improve forecast error. When we further challenge this
inference by making the two response regimes harder to distinguish, recovery
weakens, but sufficiently large models are still able to recover the
relationship, with the threshold varying by model family
(App.~\ref{app:coin-city-robustness}).

\begin{figure}[!t]
\centering
\begin{minipage}[c]{0.46\linewidth}
\includegraphics[width=\linewidth]{exp2_symbol_decoding.pdf}
\end{minipage}\hfill
\begin{minipage}[c]{0.50\linewidth}
\caption{\textbf{Larger models infer an episode-local convention, then yield
 to direct evidence.} On 88 new episodes, four larger Qwen and Llama checkpoints
(14--72B) encode the changing KIV/ZOR mapping; three smaller checkpoints
(4--8B) do not reliably encode it. The decoded state follows context before
City C has a history and follows the target relationship as local cases accumulate.}
\label{fig:coin-city-symbol-decoding}
\end{minipage}
\end{figure}

This pattern matches what we would expect if models are inferring a latent
mapping: doing so becomes harder as the underlying regimes become less
distinguishable. \emph{But when a model does recover that mapping, what has it
actually inferred?} We ask whether the episode-local KIV/ZOR relationship can be
decoded from the model's hidden states. In larger open models, the decoded
representation follows the context-selected relationship before City C has its
own history, and then shifts toward City C's observed behavior as direct evidence
accumulates (Fig.~\ref{fig:coin-city-symbol-decoding}).

But decoding still leaves open whether this representation actually matters for
the forecast. Figure~\ref{fig:coin-city-causal-patch} tests this directly by
swapping the label-state representation between otherwise matched Qwen3-14B
prompts. Before City C provides any direct observations, the intervention changes
both the analogue the model names and its numerical forecast. After four City C
cases, the same swap has no reliable effect. Across episodes, the intervention
produces a regime-aligned symmetric forecast shift of \(+.333\) before target
cases and \(-.035\) after four; the latter falls within the sign-flip null range. For Qwen3-14B, this shows that
the induced mapping helps determine the forecast when target evidence is sparse,
but gives way once direct observations reveal City C's own response relationship
(App.~\ref{app:coin-city-causal-patch}).

\begin{figure}[!t]
\vspace{-0.4em}
\centering
\includegraphics[width=\linewidth]{exp2_causal_patch.pdf}
\vspace{-0.7em}
\caption{\textbf{Changing the induced convention changes the forecast when
evidence is sparse.} In the illustrated pair, swapping the internal
representation moves the forecast from 51.3 to 54.0 and changes which reference
city the model names; reversing the swap reverses the movement. As more City C
behavior is revealed, the same label-state swap has no reliable effect. Full controls in
Table~\ref{tab:coin-city-causal-patch}.}
\label{fig:coin-city-causal-patch}
\vspace{-1.0em}
\end{figure}

\noindent\textbf{Takeaway.}
Models use context to select a predictive relationship and revise that selection
from local evidence. Arbitrary labels remove a familiar-semantic shortcut;
decoding reveals an episode-local representation, and intervention links it to
Qwen3-14B's forecast when target evidence is absent. Experiment~3 asks which parts
of this behavior training carries into new domains and response structures.
